\begin{ouverture}[grandeurs]{11}{Comparer et mesurer des aires}{Partie 2 · Modéliser et mesurer}
{Surfaces · Unités carrées · Conversions · Formules}
\objectif{Comparer des aires par superposition ou par découpage.}
\objectif{Utiliser les unités cm², dm² et m² et convertir entre elles.}
\objectif{Calculer l’aire d’un carré et d’un rectangle avec une formule.}
\prerequis{Quadrillages ; multiplication des décimaux ; périmètre.}
\end{ouverture}
\pagestyle{otMain}\essentiel\fontsize{11}{15}\selectfont
\phantomsection\label{cours-11-1}
\essSection{1}{Comparer sans mesurer}
\begin{definition}L’aire d’une surface mesure la place qu’elle occupe. Deux surfaces ont la même aire si l’on peut découper l’une et recoller ses morceaux pour recouvrir exactement l’autre.\end{definition}
\label{fig-cours-11-1-cours-decoupage}\figureGeo{cours-decoupage}

\begin{attention}Aire et périmètre sont deux grandeurs différentes : en découpant et recollant, l’aire ne change pas, alors que le périmètre peut changer.\end{attention}
\essSection{2}{Comprendre les unités carrées}
\begin{definition}1 cm$^2$ est l’aire d’un carré de 1 cm de côté ; 1 dm$^2$, celle d’un carré de 1 dm de côté ; 1 m$^2$, celle d’un carré de 1 m de côté.\end{definition}
\begin{savoirfaire}[Compter les carreaux]Sur un quadrillage de carreaux de 1 cm de côté, chaque carreau a une aire de 1 cm$^2$ ; deux demi-carreaux font 1 cm$^2$.\end{savoirfaire}
\clearpage
\phantomsection\label{cours-11-2}
\essSection{3}{Convertir entre unités d’aire voisines}
\begin{propriete}$1\ \mathrm{m}^2=1\ \mathrm m\times1\ \mathrm m=10\ \mathrm{dm}\times10\ \mathrm{dm}=100\ \mathrm{dm}^2$ et $1\ \mathrm{dm}^2=100\ \mathrm{cm}^2$. Dans l’autre sens : $1\ \mathrm{cm}^2=\frac1{100}\ \mathrm{dm}^2=0{,}01\ \mathrm{dm}^2$ et $1\ \mathrm{dm}^2=0{,}01\ \mathrm{m}^2$.\end{propriete}
\label{fig-cours-11-2-cours-dm2}\figureGeo{cours-dm2}

\begin{attention}Entre deux unités d’aire voisines, le rapport est 100, et non 10 : un carré de 10 cm de côté contient 10 rangées de 10 carreaux de 1 cm$^2$.\end{attention}
\essSection{4}{Aires du carré et du rectangle}
\begin{propriete}Aire d’un carré de côté $c$ : $A=c\times c$. Aire d’un rectangle de longueur $L$ et de largeur $\ell$ : $A=L\times\ell$, les deux longueurs étant dans la même unité.\end{propriete}
\begin{exemple}Rectangle de $3{,}5$ m sur 2 m : $A=3{,}5\times2=7\ \mathrm{m}^2$. Carré de 7 cm de côté : $A=7\times7=49\ \mathrm{cm}^2$.\end{exemple}

\clearpage\rubrique{Découverte}
\begin{blocLivre}{ex-11-1}\begin{exercice}[1]{Comparer sans mesurer}
Ranger les surfaces A, B et C selon leur aire, sans règle. Deux demi-carreaux forment un carreau.
\label{fig-ex-11-1}\figureGeo{air-1}
\end{exercice}\end{blocLivre}
\begin{blocLivre}{ex-11-2}\begin{exercice}[1]{Compter les carreaux}
Chaque carreau mesure 1 cm de côté. Quelle est l’aire de la figure ?
\label{fig-ex-11-2}\figureGeo{air-2}
\end{exercice}\end{blocLivre}
\begin{blocLivre}{ex-11-3}\begin{exercice}[1]{Choisir l’unité}
Quelle unité, parmi cm$^2$, dm$^2$ et m$^2$, convient le mieux pour l’aire : d’un timbre ; d’une page de cahier ; du sol d’une chambre ?
\end{exercice}\end{blocLivre}
\begin{blocLivre}{ex-11-4}\begin{exercice}[1]{Avec les formules}
Calculer l’aire d’un carré de 7 cm de côté, puis celle d’un rectangle de 12 m sur $4{,}5$ m.
\end{exercice}\end{blocLivre}
\begin{blocLivre}{ex-11-5}\begin{exercice}[1]{Convertir}
Convertir : 3 m$^2$ en dm$^2$ ; 5 dm$^2$ en cm$^2$ ; 400 cm$^2$ en dm$^2$ ; 250 dm$^2$ en m$^2$.
\end{exercice}\end{blocLivre}

\rubrique{Entraînement}
\begin{blocLivre}{ex-11-6}\begin{exercice}[2]{Conversions décimales}
Convertir : $0{,}7$ m$^2$ en dm$^2$ ; 45 cm$^2$ en dm$^2$ ; $3{,}2$ dm$^2$ en cm$^2$ ; 8 dm$^2$ en m$^2$.
\end{exercice}\end{blocLivre}
\begin{blocLivre}{ex-11-7}\begin{exercice}[2]{Même aire, autre périmètre}
Calculer l’aire et le périmètre de chacun des deux rectangles. Que remarque-t-on ?
\label{fig-ex-11-7}\figureGeo{air-7}
\end{exercice}\end{blocLivre}
\begin{blocLivre}{ex-11-8}\begin{exercice}[2]{Découper et recoller}
On découpe le triangle grisé et on le recolle de l’autre côté. Quelle figure obtient-on ? En déduire l’aire de la figure de départ.
\label{fig-ex-11-8}\figureGeo{air-8}
\end{exercice}\end{blocLivre}
\begin{blocLivre}{ex-11-9}\begin{exercice}[2]{Retrouver une longueur}
Un rectangle a une aire de 36 cm$^2$ et une longueur de 9 cm. Quelle est sa largeur ? Un carré a une aire de 64 m$^2$ : quelle est la longueur de son côté ?
\end{exercice}\end{blocLivre}
\begin{blocLivre}{ex-11-10}\begin{exercice}[2]{La pièce en L}
Les angles du plan sont droits. Calculer l’aire de la pièce de deux façons.
\label{fig-ex-11-10}\figureGeo{air-10}
\end{exercice}\end{blocLivre}

\rubrique{Problèmes}
\begin{blocLivre}{ex-11-11}\begin{exercice}[3]{Le carrelage}
On carrelle le sol d’une pièce rectangulaire de 4 m sur $3{,}5$ m avec des carreaux carrés de 1 dm de côté. Combien de carreaux faut-il ?
\end{exercice}\end{blocLivre}
\begin{blocLivre}{ex-11-12}\begin{exercice}[3]{La peinture}
Un mur mesure 5 m sur $2{,}5$ m ; il comporte une porte de $0{,}9$ m sur 2 m, qu’on ne peint pas. Un pot de peinture couvre 5 m$^2$. Combien de pots faut-il acheter ?
\end{exercice}\end{blocLivre}
\begin{blocLivre}{ex-11-13}\begin{exercice}[3]{Le jardin partagé}
Un jardin carré de 20 m de côté est partagé entre un potager rectangulaire de 20 m sur 8 m et une pelouse. Calculer l’aire du potager, puis celle de la pelouse. Quel pourcentage du jardin le potager occupe-t-il ?
\label{fig-ex-11-13}\figureGeo{air-13}
\end{exercice}\end{blocLivre}
\begin{blocLivre}{ex-11-14}\begin{exercice}[3]{Des timbres sur une affiche}
Des timbres rectangulaires de 3 cm sur $2{,}5$ cm recouvrent sans chevauchement une affiche de 6 dm sur 4 dm. Combien de timbres faut-il ?
\end{exercice}\end{blocLivre}
\begin{blocLivre}{ex-11-15}\begin{exercice}[3]{Vrai ou faux}
Lucie affirme : « Si on double le côté d’un carré, son aire double. » Vérifier avec un carré de 3 cm de côté.
\end{exercice}\end{blocLivre}

\clearpage
